\section{Defects4J: il benchmark di valutazione}
\label{sec:defects4j}

L'ultimo elemento di contesto necessario prima di affrontare l'architettura riguarda i \textit{dati}
su cui HybridRepair viene sviluppato e messo alla prova. Un sistema di riparazione automatica ha senso
solo se valutato su difetti \emph{reali}, riproducibili e confrontabili con quelli usati dagli altri
lavori della letteratura, per questo motivo si è scelto \textbf{Defects4J} \cite{just2014defects4j}.

\subsection{Che cos'è Defects4J?}

Defects4J è una raccolta curata di \textit{bug} reali, estratti dalla storia di sviluppo di progetti
\textit{open source} Java molto noti (tra cui JFreeChart, Google Closure Compiler, Apache Commons Lang,
Apache Commons Math e Jackson).

Per ciascun difetto, il \textit{framework} mette a disposizione un ambiente di esecuzione riproducibile
e, in particolare, quattro elementi fondamentali:

\begin{itemize}
    \item la \textbf{versione difettosa} (\textit{buggy}) del codice sorgente, ovvero il programma
    immediatamente prima della correzione;
    \item la \textbf{versione corretta} (\textit{fixed}), che differisce dalla precedente per la sola
    \textit{patch} applicata dagli sviluppatori e funge da soluzione di riferimento
    (\textit{ground truth});
    \item almeno un \textbf{test che innesca il difetto} (\textit{triggering test}), cioè un caso di
    test che fallisce sulla versione \textit{buggy} e passa su quella \textit{fixed};
    \item la \textbf{suite di test di regressione} associata, usata per verificare che la correzione
    non comprometta il comportamento preesistente.
\end{itemize}

Un aspetto centrale è l'\textbf{isolamento del difetto}: la differenza tra versione \textit{buggy} e
\textit{fixed} è minimizzata e ricondotta al solo bug in esame, depurata da modifiche estranee
(refactoring, cambi di formattazione, aggiornamenti non correlati). Questo consente di attribuire in
modo inequivocabile un eventuale successo o fallimento alla \textit{patch} generata, e non a fattori
di contorno.

\subsection{Identificazione dei difetti e sottoinsieme NPE}

Ogni difetto è identificato da una sigla nella forma \texttt{Progetto-Numero} --- ad esempio
\texttt{Chart-2}, \texttt{Closure-2} o \texttt{Codec-13} --- dove la prima parte indica il progetto di
provenienza e la seconda l'indice del bug all'interno di quel progetto. È questa la notazione che il
sistema riceve in ingresso (il \texttt{bug\_id}) e che ricorre nei capitoli successivi quando si
illustrano casi concreti di riparazione.

Poiché questa tesi si concentra sulle \textit{NullPointerException} (per le ragioni discusse nel
Capitolo precedente), non si utilizza l'intero benchmark, bensì il \textbf{sottoinsieme dei difetti
riconducibili a NPE}: bug il cui \textit{triggering test} fallisce proprio per via di una dereferenziazione
di un riferimento \texttt{null}. Questo sottoinsieme costituisce il dominio di valutazione di HybridRepair.

\subsection{Il ruolo di Defects4J in HybridRepair}

All'interno del nostro lavoro Defects4J svolge un duplice ruolo, coerente con la \textit{pipeline}
descritta nei capitoli seguenti:

\begin{itemize}
    \item come \textbf{sorgente di input}, fornisce per ogni \texttt{bug\_id} il codice sorgente
    difettoso e i test falliti, ovvero esattamente il materiale da cui parte la Fase~1 (FaultOracle)
    per la localizzazione causale;
    \item come \textbf{oracolo di valutazione}, fornisce sia i test (che certificano se la \textit{patch}
    risolve davvero il difetto senza introdurre regressioni) sia la versione \textit{fixed} di
    riferimento, utile per un confronto qualitativo tra la correzione generata e quella umana.
\end{itemize}

